\chapter{SVM Classifier Evaluation}

\section{Experimental Setup}

We evaluated the Support Vector Machine (SVM) classifier on the GSE42568 breast
cancer dataset (121 samples: 104 cancer, 17 normal) using a 5-fold stratified
cross-validation strategy to prevent data leakage. This approach follows the
dual-track methodology outlined in the project architecture:

\begin{itemize}
    \item \textbf{Path A (Baseline):} Train SVM on all preprocessed features
          (54,675 genes) with no feature selection.
    \item \textbf{Path B (Optimised):} Apply a feature selection method strictly
          inside each training fold, then train SVM on the selected 20 features.
\end{itemize}

All SVM models used a linear kernel, balanced class weights to account for the
severe class imbalance (104:17), and \texttt{max\_iter=2{,}000} to bound
worst-case training time on the high-dimensional input.

%---------------------------------------------------------------------------
\section{Path A: Baseline Track Results}
%---------------------------------------------------------------------------

\subsection{Cross-Validation Performance}

The baseline SVM trained on all 54,675 genes achieved the following results
across five folds (Table~\ref{tab:path_a}).

\begin{table}[H]
    \centering
    \caption{SVM Baseline (Path A) -- Per-Fold and Aggregate Results}
    \label{tab:path_a}
    \begin{tabular}{lccccccc}
        \toprule
        \textbf{Fold} & \textbf{Acc.} & \textbf{Prec.} & \textbf{Recall}
                      & \textbf{F1} & \textbf{MCC} & \textbf{ROC-AUC}
                      & \textbf{Spec.} \\
        \midrule
        1 & 0.8400 & 0.8400 & 1.0000 & 0.9130 & 0.0000 & -- & 0.0000 \\
        2 & 0.8750 & 0.8750 & 1.0000 & 0.9333 & 0.0000 & -- & 0.0000 \\
        3 & 0.8750 & 0.8750 & 1.0000 & 0.9333 & 0.0000 & -- & 0.0000 \\
        4 & 0.8750 & 0.8750 & 1.0000 & 0.9333 & 0.0000 & -- & 0.0000 \\
        5 & 0.8333 & 0.8333 & 1.0000 & 0.9091 & 0.0000 & -- & 0.0000 \\
        \midrule
        \textbf{Mean} & \textbf{0.8597} & \textbf{0.8597} & \textbf{1.0000}
                      & \textbf{0.9244} & \textbf{0.0000} & \textbf{0.4569}
                      & \textbf{0.0000} \\
        \textbf{Std}  & 0.0211 & 0.0211 & 0.0000
                      & 0.0123 & 0.0000 & 0.1320
                      & 0.0000 \\
        \bottomrule
    \end{tabular}
\end{table}

\subsection{Key Findings -- Baseline}

Despite a nominal accuracy of 85.97\%, the baseline results reveal a critical
failure of classification on this heavily imbalanced dataset:

\begin{itemize}
    \item \textbf{MCC = 0.0000 across all five folds.}  Matthews Correlation
          Coefficient (MCC) is zero when the classifier predicts only one class.
          Combined with Recall = 1.0000 and Specificity = 0.0000, this confirms
          that the baseline SVM predicted \emph{cancer} for every sample in every
          fold -- a majority-class collapse.

    \item \textbf{Accuracy is misleading here.}  With 104 cancer and 17 normal
          samples per fold, predicting the majority class unconditionally yields
          approximately $104/121 \approx 85.9\%$ accuracy -- exactly what is
          observed.  This figure carries no discriminative information.

    \item \textbf{ROC-AUC = 0.4569.}  A value below 0.5 indicates the model
          performs worse than random chance when the decision threshold is varied,
          confirming that the linear SVM could not find a useful decision
          boundary in the 54,675-dimensional space within the iteration limit.

    \item \textbf{Implication.}  Training on 54,675 features with a small,
          imbalanced sample ($n = 121$) places this problem firmly in the
          $p \gg n$ regime.  The feature space is so vastly over-specified that
          the margin-maximisation objective becomes ill-conditioned.  Feature
          selection is not merely beneficial here -- it is necessary for the
          classifier to function at all.
\end{itemize}

%---------------------------------------------------------------------------
\section{Path B: Optimised Track Results}
%---------------------------------------------------------------------------

\subsection{Filter Method: Welch's T-Test}

\begin{table}[H]
    \centering
    \caption{SVM with Filter (Welch's T-Test, $k=20$) -- Results}
    \label{tab:ttest}
    \begin{tabular}{lcccccc}
        \toprule
        \textbf{Fold} & \textbf{Acc.} & \textbf{Prec.} & \textbf{Recall}
                      & \textbf{F1} & \textbf{MCC} & \textbf{ROC-AUC} \\
        \midrule
        1 & 0.1600 & -- & 0.0000 & 0.0000 & 0.0000 & -- \\
        2 & 0.6250 & -- & -- & 0.7429 & 0.1917 & -- \\
        3 & 0.8750 & -- & -- & 0.9333 & 0.0000 & -- \\
        4 & 0.8750 & -- & -- & 0.9333 & 0.0000 & -- \\
        5 & 0.5833 & -- & -- & 0.6667 & 0.3780 & -- \\
        \midrule
        \textbf{Mean} & \textbf{0.6237} & \textbf{0.7357} & \textbf{0.6238}
                      & \textbf{0.6552} & \textbf{0.1139} & \textbf{0.6291} \\
        \textbf{Std}  & 0.2928 & 0.4145 & 0.4146
                      & 0.3846 & 0.1693 & 0.1764 \\
        \bottomrule
    \end{tabular}
\end{table}

\paragraph{Analysis:}
Although the mean accuracy (62.37\%) is lower than the baseline, the T-Test
method is the first approach that produces a non-zero MCC (0.1139), confirming
the model has begun to discriminate between classes.  The extreme variability
-- accuracy ranging from 16.00\% (Fold 1) to 87.50\% (Folds 3 and 4) --
reflects the instability of the Welch t-test under this class imbalance.
In a fold where the 17 normal samples are unevenly distributed, the t-test
p-value ranking shifts substantially, selecting an entirely different feature
subset.  The ROC-AUC of 0.6291 is above chance, indicating partial
discriminative power.

\subsection{Filter Method: ANOVA F-Test}

\begin{table}[H]
    \centering
    \caption{SVM with Filter (ANOVA F-Test, $k=20$) -- Results}
    \label{tab:anova}
    \begin{tabular}{lcccccc}
        \toprule
        \textbf{Fold} & \textbf{Acc.} & \textbf{Prec.} & \textbf{Recall}
                      & \textbf{F1} & \textbf{MCC} & \textbf{ROC-AUC} \\
        \midrule
        1 & 0.8000 & -- & 0.7981 & 0.8837 & 0.1746 & -- \\
        2 & 0.5833 & -- & -- & 0.7368 & $-$0.2425 & -- \\
        3 & 0.7917 & -- & -- & 0.8780 & 0.1690 & -- \\
        4 & 0.7083 & -- & -- & 0.8205 & 0.0727 & -- \\
        5 & 0.6667 & -- & -- & 0.8000 & $-$0.2000 & -- \\
        \midrule
        \textbf{Mean} & \textbf{0.7100} & \textbf{0.8552} & \textbf{0.7981}
                      & \textbf{0.8238} & $-$\textbf{0.0052} & \textbf{0.5219} \\
        \textbf{Std}  & 0.0904 & 0.0426 & 0.0914
                      & 0.0606 & 0.2019 & 0.1721 \\
        \bottomrule
    \end{tabular}
\end{table}

\paragraph{Analysis:}
ANOVA improves mean accuracy to 71.00\% and dramatically reduces fold-to-fold
variance in precision (std = 0.0426 vs 0.4145 for T-Test), demonstrating
greater stability.  However, the mean MCC is essentially zero ($-$0.0052),
indicating that despite reasonable recall (0.7981) the model consistently
achieves near-zero specificity (0.1833 mean), struggling to identify the rare
normal class.  The constant-feature warnings produced during this run
(approximately 100 zero-variance probes in each fold) reflect the dataset's
imbalance and are a known artefact of Affymetrix arrays on small normal cohorts.

\subsection{Wrapper Method: SVM-RFE}

\begin{table}[H]
    \centering
    \caption{SVM with Wrapper (SVM-RFE, $k=20$) -- Results}
    \label{tab:svmrfe}
    \begin{tabular}{lcccccc}
        \toprule
        \textbf{Fold} & \textbf{Acc.} & \textbf{Prec.} & \textbf{Recall}
                      & \textbf{F1} & \textbf{MCC} & \textbf{ROC-AUC} \\
        \midrule
        1 & 0.8400 & -- & 0.8557 & 0.9048 & 0.4048 & -- \\
        2 & 0.7917 & -- & -- & 0.8780 & 0.1690 & -- \\
        3 & 0.7917 & -- & -- & 0.8837 & $-$0.1140 & -- \\
        4 & 0.6667 & -- & -- & 0.8000 & $-$0.1939 & -- \\
        5 & 0.7500 & -- & -- & 0.8500 & 0.1000 & -- \\
        \midrule
        \textbf{Mean} & \textbf{0.7680} & \textbf{0.8721} & \textbf{0.8557}
                      & \textbf{0.8633} & \textbf{0.0732} & \textbf{0.5379} \\
        \textbf{Std}  & 0.0650 & 0.0287 & 0.0584
                      & 0.0404 & 0.2378 & 0.2560 \\
        \bottomrule
    \end{tabular}
\end{table}

\paragraph{Analysis:}
SVM-RFE achieves the highest mean accuracy among the three filter/wrapper methods
evaluated here (76.80\%) and the best mean F1-Score (0.8633).  Unlike the filter
methods, SVM-RFE accounts for linear separability in the selected feature
subspace, producing features that are jointly useful rather than individually
discriminative.  The positive mean MCC (0.0732) confirms genuine bi-class
discrimination, and Fold 1 attains an MCC of 0.4048 -- a meaningful result.
The high MCC standard deviation (0.2378) reflects fold sensitivity, but the
tight precision standard deviation (0.0287) shows consistent positive predictive
value across folds.

\subsection{Wrapper Method: Greedy Forward Selection}
\label{subsec:forward}

\begin{table}[H]
    \centering
    \caption{SVM with Wrapper (Forward Selection, $k=20$) -- Results}
    \label{tab:forward}
    \begin{tabular}{lcccccc}
        \toprule
        \textbf{Fold} & \textbf{Acc.} & \textbf{Prec.} & \textbf{Recall}
                      & \textbf{F1} & \textbf{MCC} & \textbf{ROC-AUC} \\
        \midrule
        1 & 0.7600 & -- & 0.7505 & 0.8421 & 0.4023 & -- \\
        2 & 0.6667 & -- & -- & 0.8000 & $-$0.1939 & -- \\
        3 & 0.6250 & -- & -- & 0.7692 & $-$0.2182 & -- \\
        4 & 0.6667 & -- & -- & 0.7895 & 0.0346 & -- \\
        5 & 0.6667 & -- & -- & 0.8000 & $-$0.2000 & -- \\
        \midrule
        \textbf{Mean} & \textbf{0.6770} & \textbf{0.8598} & \textbf{0.7505}
                      & \textbf{0.8002} & $-$\textbf{0.0350} & \textbf{0.4303} \\
        \textbf{Std}  & 0.0498 & 0.0541 & 0.0365
                      & 0.0266 & 0.2656 & 0.1654 \\
        \bottomrule
    \end{tabular}
\end{table}

\paragraph{Analysis:}
Forward selection achieved 67.70\% mean accuracy and a mean MCC of $-$0.035,
placing it among the weakest methods on this dataset.  The ROC-AUC of 0.4303
is below random chance, indicating poor threshold-averaged discrimination.
These results are, however, secondary to the critical computational issue
documented in Section~\ref{sec:forward_complexity}.

\subsection{Wrapper Method: Backward Elimination}

\begin{table}[H]
    \centering
    \caption{SVM with Wrapper (Backward Elimination, $k=20$) -- Results}
    \label{tab:backward}
    \begin{tabular}{lcccccc}
        \toprule
        \textbf{Fold} & \textbf{Acc.} & \textbf{Prec.} & \textbf{Recall}
                      & \textbf{F1} & \textbf{MCC} & \textbf{ROC-AUC} \\
        \midrule
        1 & 0.7200 & -- & 0.8543 & 0.8293 & 0.0546 & -- \\
        2 & 0.8333 & -- & -- & 0.9048 & 0.2381 & -- \\
        3 & 0.8750 & -- & -- & 0.9302 & 0.3419 & -- \\
        4 & 0.7917 & -- & -- & 0.8837 & $-$0.1140 & -- \\
        5 & 0.6250 & -- & -- & 0.7568 & $-$0.0410 & -- \\
        \midrule
        \textbf{Mean} & \textbf{0.7690} & \textbf{0.8702} & \textbf{0.8543}
                      & \textbf{0.8609} & \textbf{0.0959} & \textbf{0.5283} \\
        \textbf{Std}  & 0.0988 & 0.0365 & 0.1007
                      & 0.0691 & 0.1905 & 0.1103 \\
        \bottomrule
    \end{tabular}
\end{table}

\paragraph{Analysis:}
Backward elimination achieves a comparable mean accuracy (76.90\%) and F1-Score
(0.8609) to SVM-RFE, and the highest mean MCC among the wrapper methods
(0.0959).  Folds 2 and 3 produce particularly strong results (MCC 0.24 and
0.34 respectively).  However, its practical utility is severely limited by its
computational cost, detailed in Section~\ref{sec:forward_complexity}, and by
the observation (Section~4 of the feature selection report) that backward
elimination on the 500-feature pre-filtered pool tends to converge on
unannotated late-array probes rather than biologically validated genes.

\subsection{Embedded Method: LASSO / L1 Logistic Regression}

\begin{table}[H]
    \centering
    \caption{SVM with Embedded Feature Selection (LASSO, $k \leq 20$) -- Results}
    \label{tab:lasso}
    \begin{tabular}{lcccccc}
        \toprule
        \textbf{Fold} & \textbf{Acc.} & \textbf{Prec.} & \textbf{Recall}
                      & \textbf{F1} & \textbf{MCC} & \textbf{$k$ used} \\
        \midrule
        1 & 0.7600 & -- & 0.9710 & 0.8636 & $-$0.1287 & 18 \\
        2 & 0.8750 & -- & -- & 0.9333 & 0.0000 & 20 \\
        3 & 0.8750 & -- & -- & 0.9333 & 0.0000 & 20 \\
        4 & 0.8750 & -- & -- & 0.9333 & 0.0000 & 20 \\
        5 & 0.7917 & -- & -- & 0.8837 & $-$0.0933 & 18 \\
        \midrule
        \textbf{Mean} & \textbf{0.8353} & \textbf{0.8554} & \textbf{0.9710}
                      & \textbf{0.9095} & $-$\textbf{0.0444} & -- \\
        \textbf{Std}  & 0.0555 & 0.0268 & 0.0429
                      & 0.0334 & 0.0621 & -- \\
        \bottomrule
    \end{tabular}
\end{table}

\paragraph{Analysis:}
LASSO achieves the highest mean accuracy of all Path B methods (83.53\%) and
the highest mean F1-Score (0.9095), driven by an exceptional mean recall of
0.9710 -- the model successfully identifies 97\% of cancer samples.  The very
low MCC standard deviation (0.0621) makes it the most \emph{consistent} method
across folds.  However, the mean MCC of $-$0.0444 and specificity of 0.0000
reveal the same majority-class bias as the baseline: LASSO selects features
that collectively predict cancer but cannot reliably distinguish the rare normal
class.  Notably, LASSO produced only 18 non-zero-coefficient features in Folds
1 and 5, reflecting the sparsity of the L1 penalty at $C = 0.1$.

%---------------------------------------------------------------------------
\section{Comparative Summary}
%---------------------------------------------------------------------------

Table~\ref{tab:comparison} presents the cross-method comparison using all
metrics reported in the experimental output.

\begin{table}[H]
    \centering
    \caption{SVM Performance Comparison Across All Methods (GSE42568)}
    \label{tab:comparison}
    \small
    \begin{tabular}{lcccccc}
        \toprule
        \textbf{Method} & \textbf{Accuracy} & \textbf{F1} & \textbf{MCC}
                        & \textbf{Recall} & \textbf{Specificity}
                        & \textbf{Features} \\
        \midrule
        Path A: Baseline
            & $0.8597 \pm 0.0211$ & $0.9244 \pm 0.0123$ & $0.0000 \pm 0.0000$
            & $1.0000 \pm 0.0000$ & $0.0000 \pm 0.0000$ & 54,675 \\
        Path B: T-Test
            & $0.6237 \pm 0.2928$ & $0.6552 \pm 0.3846$ & $0.1139 \pm 0.1693$
            & $0.6238 \pm 0.4146$ & $0.5333 \pm 0.5055$ & 20 \\
        Path B: ANOVA
            & $0.7100 \pm 0.0904$ & $0.8238 \pm 0.0606$ & $-0.0052 \pm 0.2019$
            & $0.7981 \pm 0.0914$ & $0.1833 \pm 0.1708$ & 20 \\
        Path B: SVM-RFE
            & $0.7680 \pm 0.0650$ & $0.8633 \pm 0.0404$ & $0.0732 \pm 0.2378$
            & $0.8557 \pm 0.0584$ & $0.2167 \pm 0.2173$ & 20 \\
        Path B: Fwd Selection
            & $0.6770 \pm 0.0498$ & $0.8002 \pm 0.0266$ & $-0.0350 \pm 0.2656$
            & $0.7505 \pm 0.0365$ & $0.2167 \pm 0.3312$ & 20 \\
        Path B: Bwd Elimination
            & $0.7690 \pm 0.0988$ & $0.8609 \pm 0.0691$ & $0.0959 \pm 0.1905$
            & $0.8543 \pm 0.1007$ & $0.2333 \pm 0.1369$ & 20 \\
        Path B: LASSO
            & $0.8353 \pm 0.0555$ & $0.9095 \pm 0.0334$ & $-0.0444 \pm 0.0621$
            & $0.9710 \pm 0.0429$ & $0.0000 \pm 0.0000$ & $\leq$20 \\
        \bottomrule
    \end{tabular}
\end{table}

%---------------------------------------------------------------------------
\section{The Computational Cost of Greedy Forward Selection}
\label{sec:forward_complexity}
%---------------------------------------------------------------------------

\subsection{Mathematical Complexity}

Greedy forward selection is fundamentally an exhaustive search algorithm.
At each step $i$, the method trains a separate SVM model for every remaining
candidate feature and selects the one that maximises training accuracy.  For a
dataset with $N$ total features and a target subset of $k$ features, the total
number of SVM model fits required per cross-validation fold is:

\begin{equation}
    \text{Models per fold} = \sum_{i=0}^{k-1}(N - i)
    = kN - \frac{k(k-1)}{2}
    \label{eq:fwd_complexity}
\end{equation}

For this experiment ($N = 54{,}675$, $k = 20$):

\begin{equation}
    \text{Models per fold}
        = 20 \times 54{,}675 - \frac{20 \times 19}{2}
        = 1{,}093{,}500 - 190
        \approx 1.09 \times 10^6
\end{equation}

Across the five cross-validation folds, the pipeline must therefore train and
evaluate approximately \textbf{5.46 million distinct SVM models}.  With 96
training samples per fold on a 54,675-dimensional feature space, each model fit
is itself non-trivial.

\subsection{Observed Runtime Evidence from Training Logs}

The scale of this computational burden is directly visible in the training log.
Table~\ref{tab:fwd_timing} shows the wall-clock time between consecutive steps
for the first fold alone.

\begin{table}[H]
    \centering
    \caption{Forward Selection Timing -- Fold 1 of 5 (from training log)}
    \label{tab:fwd_timing}
    \begin{tabular}{ccc}
        \toprule
        \textbf{Step Transition} & \textbf{Clock Time} & \textbf{Elapsed} \\
        \midrule
        Step 1 (54,675 models)  & 03:34 $\to$ 03:36 & $\approx$ 2 min \\
        Step 8 $\to$ Step 9     & 04:00 $\to$ 05:05 & $\approx$ 65 min \\
        Step 12 $\to$ Step 13   & 05:12 $\to$ 09:02 & $\approx$ 230 min \\
        \bottomrule
    \end{tabular}
    \bigskip \\
    \small\textit{Note: Step 1 is fast because each SVM fits a single feature
    (trivially separable or not). As the selected set grows, each candidate
    evaluation fits a progressively wider matrix, compounding the cost.}
\end{table}

The non-linear growth in step duration is a direct consequence of the SVM's
$\mathcal{O}(n^2)$ to $\mathcal{O}(n^3)$ training complexity with respect to
sample count interacting with an ever-widening feature matrix.  The total
wall-clock time for the full forward selection run across all five folds was
approximately \textbf{16 hours} (03:34 on 12 June 2026 to approximately 19:30,
with backward elimination beginning at 17:30).

\subsection{Why the Accuracy Did Not Improve}

The 16-hour runtime produced a mean accuracy of only 67.70\% and a mean MCC of
$-$0.035 -- worse than ANOVA (71.00\%) and SVM-RFE (76.80\%), both of which
completed in under one second per fold.  Three factors explain this paradox:

\begin{enumerate}
    \item \textbf{Majority-class anchoring.}  At Step 1, any feature that
          separates the 104 cancer samples as a bloc scores approximately 0.875
          accuracy (the majority-class ratio), regardless of whether it has any
          clinical relevance.  The algorithm locks onto this score and spends
          subsequent steps adding features that maintain it rather than probing
          for normal-class discrimination.

    \item \textbf{Greedy non-optimality.}  Forward selection is a locally
          greedy heuristic.  The feature chosen at Step 1 conditions all
          subsequent choices.  If the globally optimal 20-feature subset requires
          two features that individually score below the majority-class ceiling,
          neither will ever be selected.

    \item \textbf{Interaction blindness at first step.}  On an imbalanced
          dataset, the feature that maximises single-feature training accuracy is
          typically one that happens to be expressed uniformly in cancer samples.
          The signal that distinguishes normal from cancer often resides in
          combinations of moderately expressed genes -- exactly the space that
          greedy search does not explore.
\end{enumerate}

\subsection{Practical Recommendation}

For genomic datasets with $N > 10{,}000$ features and $n < 500$ samples,
forward and backward selection should \textbf{not} be applied as in-fold
selectors.  The recommended alternative is to use the pre-computed feature
selection results produced by a dedicated feature selection script (run once
on the full dataset before cross-validation begins), caching the selected
indices and applying them consistently across folds.  This approach, implemented
in the updated \texttt{svm\_classifier.py}, reduces the effective per-fold
computation to a single matrix slice and eliminates the $10^6$-model search
entirely.

%---------------------------------------------------------------------------
\section{Key Insights}
%---------------------------------------------------------------------------

\subsection{Feature Selection Restores Classification Ability}

The most important result of this chapter is qualitative rather than
quantitative: \textbf{without feature selection, the SVM fails completely on
this dataset} (MCC = 0.0000).  Every Path B method that produced a non-zero
MCC -- even imperfect ones -- is superior to the baseline in terms of genuine
discriminative power.  The lesson is that on severely imbalanced,
high-dimensional datasets, feature selection is a prerequisite for meaningful
classification, not a post-hoc optimisation.

\subsection{Best Performers by Metric}

\begin{itemize}
    \item \textbf{Highest accuracy and consistency:} LASSO
          ($0.8353 \pm 0.0555$) with the lowest F1 standard deviation (0.0334),
          making it the most reliable method fold-to-fold.
    \item \textbf{Highest genuine discrimination (MCC):} Backward Elimination
          ($0.0959 \pm 0.1905$), though at prohibitive computational cost.
    \item \textbf{Best balance of speed and discrimination:} SVM-RFE
          ($\text{MCC} = 0.0732$, F1 = 0.8633, completes in under 1 second per
          fold with ANOVA pre-filtering to 500 features).
    \item \textbf{Best recall:} LASSO (0.9710), identifying 97.1\% of cancer
          cases -- clinically important for a screening application where
          missed positives are the primary risk.
\end{itemize}

\subsection{The Imbalance Problem Persists}

Despite feature selection, no method achieves a specificity above 0.53.
All methods with high recall exhibit near-zero specificity, and vice versa,
reflecting the fundamental difficulty of learning a boundary for 17 normal
samples in a 20-dimensional space.  Future work should consider:

\begin{itemize}
    \item \textbf{SMOTE or other oversampling} applied to the training fold
          before SVM fitting to synthesise additional normal samples.
    \item \textbf{Cost-sensitive SVM} with an asymmetric penalty that
          amplifies the cost of misclassifying normal samples beyond the
          current \texttt{class\_weight=`balanced'} setting.
    \item \textbf{Larger $k$} (e.g.\ $k = 50$ or $k = 100$) to provide the
          SVM with more geometric freedom when separating a minority class of
          only 17 samples.
\end{itemize}

\subsection{Computational Efficiency Summary}

\begin{table}[H]
    \centering
    \caption{Approximate Computational Cost per Cross-Validation Run}
    \label{tab:cost}
    \begin{tabular}{lcc}
        \toprule
        \textbf{Method} & \textbf{Total Wall Time (5 folds)}
                        & \textbf{SVM Models Trained} \\
        \midrule
        Baseline (no selection) & $\approx$25 sec  & 5 \\
        T-Test                  & $<$5 sec         & 5 + 5$\times$\text{selector} \\
        ANOVA                   & $<$5 sec         & 5 + 5$\times$\text{selector} \\
        SVM-RFE (pre-filtered)  & $\approx$2 sec   & 5 + 5$\times$\text{RFE} \\
        LASSO                   & $\approx$15 sec  & 5 + 5$\times$\text{LR} \\
        Forward Selection       & $\approx$16 hrs  & $\approx$5.46 million \\
        Backward Elimination    & $\approx$2.5 hrs & $\approx$0.9 million \\
        \bottomrule
    \end{tabular}
\end{table}

The six-order-of-magnitude difference in model evaluations between SVM-RFE and
Forward Selection, with the greedy search producing inferior results, provides a
compelling empirical argument for pre-filter-based wrapper methods on
high-dimensional genomic data.

%---------------------------------------------------------------------------
\section{Discussion}
%---------------------------------------------------------------------------

\subsection{Why the Baseline SVM Collapses on GSE42568}

The GSE42568 dataset presents the SVM with two compounding challenges.  First,
the class imbalance ratio (104:17 $\approx$ 6:1) means that the
\texttt{balanced} class weight setting must multiply the normal-class gradient
contribution by a factor of approximately 6 to equalise the loss surface --
but with only 17 normal training samples (approximately 14 per fold), there
is insufficient geometric information to define a robust separating hyperplane
in 54,675 dimensions.  Second, the vast majority of the 54,675 features are
irrelevant noise.  The SVM's margin-maximisation objective becomes degenerate
when the signal-to-noise ratio in the feature space approaches zero, and the
solver converges to the trivial solution of predicting the majority class.

\subsection{Why LASSO Has High Recall but Zero Specificity}

LASSO's L1 penalty with $C = 0.1$ selects features whose expression
differences are large in magnitude across the full training set.  On this
dataset, those are overwhelmingly cancer-upregulated genes -- features that
fire strongly in cancer but are also moderately present in normal tissue.
The resulting SVM decision boundary sits such that nearly all samples are
classified as cancer, maximising recall but collapsing specificity.  Increasing
$C$ (reducing regularisation) or using cross-validated $C$ selection would
allow the LASSO to retain features with smaller but more specific effects on
the normal class.

\subsection{SVM-RFE as the Practical Optimum}

Among all methods evaluated, SVM-RFE offers the most favourable
accuracy--MCC--computational-cost trade-off.  The ANOVA pre-filter reduces the
candidate space from 54,675 to 500 features before RFE begins, making the RFE
computationally feasible while preserving the highest-signal probes.  The
resulting feature set -- dominated by the consensus genes identified in the
feature selection report (SFTPB, AGER, NKX2-1, NAPSA, SCGB3A2, TMEM100) --
is both biologically interpretable and discriminatively useful, producing the
highest positive MCC at a tractable runtime.